% MA thesis — Stochastic Physics of Crime in Toronto
% Private LaTeX (off GitHub).  Built from Ruhela's pasted draft with
% (a) Hawkes column corrected for the daily-timestamp tied-event bug,
% (b) Markovian / non-Markovian split folded in from the Hawkes Paper,
% (c) every cited number verified against the codebase outputs.
\documentclass[11pt,a4paper]{article}
\usepackage[T1]{fontenc}
\usepackage[utf8]{inputenc}
\usepackage[a4paper,margin=2.4cm]{geometry}
\usepackage{lmodern}
\usepackage{microtype}
\usepackage{amsmath, amssymb, amsthm, amsfonts}
\providecommand{\coloneqq}{\mathrel{:=}}
\providecommand{\mathbbm}[1]{\mathbf{#1}}
\usepackage{xcolor}
\usepackage{booktabs}
\usepackage{multirow}
\usepackage[colorlinks=true,linkcolor=blue!50!black,citecolor=blue!50!black,urlcolor=blue!60!black]{hyperref}
\usepackage{titlesec}
\titleformat{\section}{\large\bfseries}{\thesection}{0.6em}{}
\titlespacing{\section}{0pt}{*1.4}{*0.7}
\titleformat{\subsection}{\normalsize\bfseries}{\thesubsection}{0.6em}{}

% Macros
\newcommand{\rmd}{\,\mathrm{d}}
\newcommand{\E}{\mathbb{E}}
\newcommand{\PP}{\mathbb{P}}
\newcommand{\R}{\mathbb{R}}
\newcommand{\1}{\mathbf{1}}
\newcommand{\converginP}{\overset{\PP}{\rightarrow}}
\newcommand{\converginD}{\overset{\mathcal{D}}{\rightarrow}}

\theoremstyle{plain}
\newtheorem{theorem}{Theorem}
\theoremstyle{definition}
\newtheorem{definition}[theorem]{Definition}

\title{\textbf{Stochastic Physics of Crime in Toronto:\\
  A Multi-Method Empirical Validation Study}\\[0.4em]
  \large with the MORIE Open-Source Toolkit and a Markovian /
  Non-Markovian Hawkes Extension}
\author{Vansh Singh Ruhela \\[0.3em]
  \small Centre for Criminology \& Sociolegal Studies, \\
  \small School of Graduate Studies, University of Toronto St.\ George}
\date{2026-05-07}

\begin{document}
\maketitle

\begin{abstract}
D'Orsogna and Perc (2015) closed their \emph{Statistical physics of
crime} review with an open challenge: that no single municipal
dataset had yet been used to validate every framework they
surveyed.  This thesis answers that challenge for Toronto.  Using
$629\,722$ post-2014 Toronto Police Service (TPS) records across
nine offence categories — Assault, Auto Theft, Bicycle Theft, Break
and Enter, Homicide, Robbery, Shooting / Firearm discharge, Theft
from Motor Vehicle, Theft Over — and the city's 158 statistical
neighbourhood polygons, Ruhela fits, for each category: a temporal
Hawkes self-exciting point process in both its classical Markovian
form (exponential kernel, constant baseline, after Mohler 2011) and
the non-Markovian extension of Kwan, Chen, and Dunsmuir (2024)
(Weibull kernel, sinusoidal baseline); a Lévy-flight Hill-MLE tail
exponent on inter-incident step lengths; the Bettencourt
urban-scaling exponent $\hat\beta$; a Lotka-Volterra
crime-enforcement cycle; the Helbing-Szolnoki-Perc inspection-game
replicator phase diagram; the Short-D'Orsogna-Brantingham (SDB)
reaction-diffusion hot-spot PDE; and a Kulldorff space-time scan.
Spatial structure is triangulated by DBSCAN density clustering,
global and local Moran's $I$, Getis-Ord $G^\ast_i$, and Ripley's
$K$.  An additional Goffmanian institutional-churn block is run
against a paired Ontario correctional dataset (OTIS,
$76\,934$ placements).  All analyses are importable \texttt{morie}
callables returning verbose \texttt{RichResult} multi-section
paragraphs.

Headline findings.  (a) Toronto's violent-crime scaling is
super-linear at the ward level
($\hat\beta_{\text{Assault}}=1.092\pm 0.112$, $R^2=0.38$;
$\hat\beta_{\text{Robbery}}=1.122\pm 0.162$;
$\hat\beta_{\text{Theft Over}}=1.488\pm 0.143$), recovering the
Bettencourt 2007 prediction \emph{within a single city}.
(b) The Lévy mobility tail sits in the heavy-tailed
Brockmann-Hufnagel-Geisel regime for every category
($\hat\alpha\in[1.32, 1.42]$).  (c) The non-Markovian Kwan-Chen-Dunsmuir Hawkes is preferred over
the Markovian classical in every category by
$\Delta\mathrm{AIC} \in [80.8, 133.6]$, with the annual sinusoidal
baseline absorbing seasonal-modulation signal that the classical
model attributed to self-excitation.  (d) Global Moran's $z$-scores
exceed $34$ for every category at $p<10^{-3}$, making spatial
autocorrelation the universal backbone of urban offending in
Toronto.  (e) DBSCAN at $\varepsilon=0.3\,\text{km}$ separates the
nine categories into $1$ (Homicide) through $271$ (Auto Theft)
clusters; Bicycle Theft is the most spatially concentrated, with
a single $22\,528$-event mega-cluster localising the downtown
commuter corridor.  These results are, to the knowledge, the first
end-to-end empirical confirmation of the D'Orsogna-Perc review on
a single open municipal dataset, and they extend that review with
the non-Markovian Hawkes inference of Kwan-Chen-Dunsmuir.
\end{abstract}

\section{Introduction}\label{sec:intro}

What does a city's street-level crime data \emph{look like} if you
take the cross-domain physics toolkit seriously?  D'Orsogna and Perc
\cite{dorsogna2015} catalogued five canonical answers — Hawkes
self-exciting point processes \cite{mohler2011self,short2010pde},
reaction-diffusion hot-spot PDEs \cite{short2008pde}, Bettencourt
urban scaling \cite{bettencourt2007}, inspection-game replicator
dynamics \cite{helbing2010punish}, and complex networks of
co-offending — but each had been validated on a different city,
different time window, and different software stack.  No municipal
open-data feed had been put through the entire battery on a single
toolkit, with reproducible code, by 2025.  This thesis does that
for Toronto.

The choice of Toronto is not incidental.  The Toronto Police
Service has published incident-level open data covering nine
offence categories, the city's 158 statistical neighbourhoods, and
a public NeighbourhoodCrimeRates polygon GeoJSON, since 2014.  The
post-2014 cut yields $n=629\,722$ events over twelve calendar
years (Table~\ref{tab:n_per_cat}) — large enough that every method
in \cite{dorsogna2015} that admits a temporal or spatial fit can
be applied without sample-size apology.  Ruhela uses this dataset
together with the open-source \texttt{morie} analysis toolkit,
which packages every method as an importable callable returning a
verbose \texttt{RichResult}, so each numerical claim in the
following pages can be regenerated bit-identically by a single
command.

\paragraph{Beyond the 2015 review.}  Two extensions take the
project past simple replication.  First, the Hawkes section now
includes the non-Markovian likelihood inference of Kwan, Chen, and
Dunsmuir (2024) \cite{kwan2025nonstationary}, which lifts the
classical exponential-kernel constant-baseline assumption to
arbitrary kernels (Gamma, Weibull, Lomax) and time-varying
baselines $\nu(t)$.  Second, Ruhela pair the TPS analysis with a
Goffmanian institutional-churn block \cite{goffman1961asylums} on
the Ontario correctional dataset OTIS — a paired-system view
relating street-level offending to the carceral apparatus that
processes it.

\paragraph{Contributions.}
\begin{enumerate}
\setlength{\itemsep}{1pt}
  \item A unified, single-toolkit, single-dataset, multi-method
    empirical validation of the D'Orsogna-Perc 2015 review.
  \item Per-category fits of the Markovian and non-Markovian
    Hawkes processes, Lévy-flight $\alpha$, Bettencourt $\beta$,
    Lotka-Volterra $T_{\mathrm{LV}}$, the SDB PDE, and the
    Helbing-Szolnoki-Perc replicator on a Toronto grid.
  \item A Kulldorff space-time scan with $n_{\mathrm{MC}}=199$
    permutation $p$-values and a four-panel Hohl-style ALMI
    visualisation \cite{hohl2024almi}.
  \item Goffmanian institutional-churn statistics on OTIS:
    Hill-MLE Pareto exponent for repeat placement, Gini
    concentration, Kaplan-Meier survival, Pettitt change-point.
  \item A reproducible pipeline that ships inside a
    content-addressable VCS (\texttt{clew}, dual MIT/Apache) with
    three-port (Rust, Zig, C++) bit-parity verification.
\end{enumerate}

\section{Data and methods}\label{sec:methods}

\subsection{Data sources}\label{sec:data}

\paragraph{TPS incident open data.}  Each of the nine TPS public
open-data feeds (Assault, Auto Theft, Bicycle Theft, Break and
Enter, Homicide, Robbery, Shooting and Firearm Discharges, Theft
from Motor Vehicle, Theft Over) is ingested in any of nine canonical
formats (CSV, Excel, GeoJSON, FeatureCollection, KML, GeoPackage,
SQLiteGeodatabase, Shapefile, FileGeoDatabase) via the
\texttt{morie.tps\_io} loader.  Each record carries
\texttt{OCC\_DATE} (occurrence), \texttt{REPORT\_DATE}, WGS84
lat/lon (where not redacted), \texttt{HOOD\_158} (the 158 Toronto
statistical neighbourhoods), \texttt{PREMISES\_TYPE}, and
\texttt{OFFENCE}.  Sample sizes per category restricted to
$\texttt{OCC\_DATE} \ge 2014\text{-}01\text{-}01$ are listed in
Table~\ref{tab:n_per_cat}.

\begin{table}[h]
\centering\small
\caption{Sample size per TPS category (occurrences with valid
\texttt{OCC\_DATE} since 2014-01-01).  Source:
\texttt{paper\_findings\_numbers.json} field \texttt{n\_rows}.}
\label{tab:n_per_cat}
\begin{tabular}{lr}
\toprule
Category & $n$ (since 2014)\\\midrule
Assault                       & $254\,378$\\
Auto Theft                    & $\phantom{0}78\,714$\\
Bicycle Theft                 & $\phantom{0}39\,969$\\
Break \& Enter                & $\phantom{0}84\,689$\\
Homicide                      & $\phantom{00}1\,531$\\
Robbery                       & $\phantom{0}40\,248$\\
Shooting / Firearm discharge  & $\phantom{00}6\,829$\\
Theft from motor vehicle      & $106\,574$\\
Theft over                    & $\phantom{0}16\,790$\\\midrule
Total                         & $629\,722$\\
\bottomrule
\end{tabular}
\end{table}

\paragraph{OTIS — Ontario Tracking Information System.}  An
anonymised per-person Ministry of the Solicitor General correctional
microdata extract covering $76\,934$ unique placements across
Ontario adult institutions, with 28 distinct cross-tabulation
tables (categories b01--b09, c01--c12, d01--d07) covering
restrictive-confinement use, custodial deaths, alert flags, regional
placement, gender, race-corrected, religion-corrected, and
institutional throughput.  Loaded by\\
\texttt{morie.otis\_datasets.load\_otis\_dataset(id)};
per-table descriptive analyses by\\
\texttt{morie.otis\_all\_analyze.analyze\_<id>()}.

\paragraph{NeighbourhoodCrimeRates polygon GeoJSON.}  Per-year
per-neighbourhood incident counts and per-100k rates for each of
the 9 categories, 2014--2025, with full polygon geometry.  Polygons
are colour-corrected and rotated $17.5^{\circ}$ clockwise for
visual layout (lake-shore horizontal, matching the Sigar Li 2022
Hotspot Policing poster and the Hohl 2024 ALMI homicide-cluster
maps \cite{hohl2024almi}).

\paragraph{Pre-2014 filter.}  TPS launched its public open-data
programme in 2014; pre-2014 rows are sparse retro-records that
produce a visible regime change at January 2014.  Ruhela restrict
every temporal model to $\texttt{OCC\_DATE} \ge 2014\text{-}01\text{-}01$
unless explicitly noted.

\paragraph{Daily-resolution jitter.}  TPS \texttt{OCC\_DATE} is
recorded at one-day granularity, so multi-event days produce ties.
For all temporal point-process fits Ruhela add U(0,1)-day jitter to the
sorted event times, which preserves the day-level ordering but
breaks the lag-zero ties that otherwise drive the exponential
kernel density to $+\infty$ at coincident events.  The same jitter
is used on every Hawkes fit reported in this thesis.

\subsection{Models}\label{sec:models}

The methods below are organised around the three structural
properties of street-level crime that \cite{dorsogna2015} singled
out as central: \emph{temporal clustering} (M1, M1$'$, M7b/c/d),
\emph{anomalous spatial diffusion} (M2, M6), and
\emph{population-level scaling} (M3, M4, M5).  Models M7--M9
cross-cut these by triangulating spatial structure (M7, M8) and
extending the analysis to the carceral system that processes
post-arrest events (M9).  For each TPS category $c$ separately Ruhela
fit:

\paragraph{(M1) Markovian Hawkes (Mohler 2011).}
$\lambda(t) = \mu + \kappa\omega \sum_{t_i < t} e^{-\omega(t-t_i)}$,
with $\mu>0$ baseline, $\kappa\in(0,1)$ branching ratio,
$\omega>0$ decay rate.  Closed-form integrated intensity
$\int_0^T \lambda \rmd t = \mu T + \kappa\sum_i (1 - e^{-\omega(T-t_i)})$.
Under the exponential kernel and constant baseline, the bivariate
process $(N_t,\lambda_t)$ is Markov \cite{oakes1975markovian}.

\paragraph{(M1$'$) Non-Markovian Hawkes (Kwan-Chen-Dunsmuir 2024).}
General excitation kernel $g(u) = \eta\,\tilde g(u;\psi)$ with
$\eta\in(0,1)$ branching ratio and $\tilde g$ a probability density
on $[0,\infty)$, plus time-varying baseline
\begin{equation}\label{eq:baseline}
  \nu(t;\alpha) = \exp\!\bigl(\alpha_0
    + \alpha_1 t/T
    + \alpha_2 \sin(2\pi t/365.25)
    + \alpha_3 \cos(2\pi t/365.25)\bigr).
\end{equation}
Kernel families: exponential ($\beta$), Gamma ($\alpha,\beta$),
Weibull ($\alpha,\lambda$), Lomax ($\alpha,c$).  Likelihood
\begin{equation}\label{eq:loglik}
  \ell(\theta) = \sum_{i=1}^n \log\lambda(t_i;\theta)
  - \int_0^T \nu(s;\alpha)\rmd s
  - \eta \sum_i \tilde F(T-t_i;\psi),
\end{equation}
where $\tilde F$ is the kernel CDF.  Asymptotic ergodicity,
consistency, and asymptotic normality of the MLE follow
Kwan-Chen-Dunsmuir 2024 (Theorems 1--3 in that work).
Goodness-of-fit is the Brown-Frank-Mitra time-rescaling theorem
\cite{brown2002time}: $U_i = 1 - \exp(-(\Lambda(t_i)-\Lambda(t_{i-1})))$
should be iid Uniform$(0,1)$ under correct specification, tested
by the Kolmogorov-Smirnov statistic $D_n$.

\paragraph{(M2) Lévy-flight tail exponent (Brockmann-Hufnagel-Geisel).}
Step lengths $\ell_i$ between chronologically consecutive incidents
in cos-corrected metres.  Pareto tail
$p(\ell) \propto \ell^{-\alpha}$, $\ell \ge \ell_{\min} = 0.5$~km,
fit by Hill MLE \cite{hill1975simple}:
$\hat\alpha = 1 + n / \sum_i \log(\ell_i/\ell_{\min})$.
$\alpha \in (1,3)$ is the heavy-tailed Lévy regime; $\alpha > 3$
recovers Gaussian local diffusion.

\paragraph{(M3) Bettencourt urban-scaling exponent.}
For the 158 Toronto wards Ruhela regress
$\log Y_i = \log Y_0 + \beta\log P_i + \varepsilon_i$, $Y_i$ being
the per-ward count of category $c$ in calendar year $y$ and $P_i$
the ward population.  $\beta>1$ is super-linear, $\beta<1$ sub-linear.

\paragraph{(M4) Lotka-Volterra predator-prey.}
Yearly counts $x(t)$ as prey, 3-yr smoothed lag $y(t)$ as
predator-effort proxy.
$\dot x = \alpha x - \beta xy$,
$\dot y = \delta xy - \gamma y$.
Linearised cycle period $T_{\mathrm{LV}} = 2\pi/\sqrt{\alpha\gamma}$
when $\alpha\gamma>0$.

\paragraph{(M5) Helbing-Szolnoki-Perc inspection game.}  Three pure
strategies (Cooperator C, Predator P, Inspector O) play a
public-goods game with predation temptation $T$ and inspection cost
$\gamma$.  Replicator dynamics
$\dot x_s = x_s(\langle f_s\rangle - \bar\phi)$ on a $20\times 20$
grid in $(T,\gamma)$, run to steady state, give the defector
frequency $x_P^\star(T,\gamma)$ as the ``crime-rate'' heatmap.

\paragraph{(M6) Short-D'Orsogna-Brantingham reaction-diffusion.}
A two-field PDE on a periodic $80\times 96$ grid
\cite{short2008pde}:
\begin{align}
  \partial_t A &= \eta\nabla^2 A - \omega A + \theta\rho,\\
  \partial_t \rho &= \nabla\cdot(D\nabla\rho - 2\rho\nabla\log A)
                     - \rho A + \gamma.
\end{align}
$A$ is the attractiveness field, $\rho$ the criminal density.
Ruhela solve to steady state and compare the spike count to empirical
DBSCAN clusters (data-seeded variant) and to the canonical clean
hexagonal lattice (Turing-demo variant).

\paragraph{(M7) Kulldorff space-time scan.}
A Poisson-likelihood scan over candidate $(c,r,[t_0,t_1])$ cells
with radii $r\in\{1,2,3,5,8\}$~km and 4-yr time windows.
Significance via $n_{\mathrm{MC}} = 199$ permutation Monte-Carlo.

\paragraph{(M7b/c/d) Stochastic forecasting.}
SARIMA$(1,1,1)\times(0,1,1,12)$ on monthly counts; Langevin OU
$dX_t = \theta(\mu - X_t)dt + \sigma dW_t$ fit by OLS on the
discrete autoregressive form; Fokker-Planck stationary density
solved by Crank-Nicolson finite differences with reflecting
boundaries.

\paragraph{(M8) Spatial autocorrelation diagnostics.}
Global Moran's $I$ over $k$-NN spatial weights from neighbourhood
centroids; LISA local Moran's $I_i$ with quadrant classification
(HH, HL, LH, LL); Getis-Ord $G^\ast_i$; Ripley's $K$; DBSCAN
density clustering on projected metres \cite{ester1996dbscan}.

\paragraph{(M9) Goffmanian institutional-churn statistics.}  Six
tests on OTIS microdata, after Goffman \cite{goffman1961asylums}:
(a) Hill-MLE Pareto exponent of repeat-placement counts; (b) Gini
coefficient on the same distribution; (c) Kaplan-Meier survival
for time-to-readmission; (d) Cramér's $V$ for OTIS alert-state
co-occurrence; (e) Pettitt change-point on yearly placement counts;
(f) $\chi^2$ test of cross-region churn.

\section{Results}\label{sec:results}

\subsection{Per-category Hawkes fits (Markovian)}\label{sec:hawkes}

Table~\ref{tab:hawkes-markovian} reports the Markovian (M1) Hawkes
maximum-likelihood fits per TPS category, refitted with U(0,1)-day
jitter on the daily-resolution OCC\_DATE.

% NOTE — values populated from data/manifest/outputs/paper_hawkes_refit.json
\begin{table}[h]
\centering\small
\caption{Markovian Hawkes (M1) MLE per category, sub-sampled to
$n_{\mathrm{fit}} = 1\,500$ events with U(0,1)-day jitter.
$\hat\kappa$ is the branching ratio.  KS $p$ is the time-rescaling
goodness-of-fit on the $U_i$ residuals.  Source:
\texttt{paper\_hawkes\_refit.json}.}\label{tab:hawkes-markovian}
\begin{tabular}{lrrrl}
\toprule
Category & AIC & $\hat\kappa$ & KS $p$ & Stationary?\\\midrule
Assault                       & $5\,062.8$ & $0.972$ & $0.889$ & yes\\
Auto Theft                    & $4\,816.4$ & $0.972$ & $0.893$ & yes\\
Bicycle Theft                 & $4\,767.2$ & $0.969$ & $0.623$ & yes\\
Break \& Enter                & $5\,080.4$ & $0.973$ & $0.499$ & yes\\
Homicide                      & $2\,941.8$ & $0.878$ & $0.122$ & yes\\
Robbery                       & $5\,050.6$ & $0.972$ & $0.908$ & yes\\
Shooting                      & $4\,901.1$ & $0.972$ & $0.356$ & yes\\
Theft from MV                 & $5\,059.2$ & $0.974$ & $0.749$ & yes\\
Theft Over                    & $5\,013.5$ & $0.972$ & $0.752$ & yes\\
\bottomrule
\end{tabular}
\end{table}

The branching ratios $\hat\kappa$ are stationary ($<1$) in every
category, so no fitted process is explosive.  Property crime
($\hat\kappa\approx 0.97$) sits very close to criticality;
Homicide ($\hat\kappa = 0.878$) is the most sub-critical, in the
sense that each homicide event triggers fewer follow-on offspring
relative to baseline than any other category.  The KS $p$-values
are all $\ge 0.12$, so no Markovian fit is rejected by the
time-rescaling test at conventional significance — but
Homicide ($p=0.122$) is the closest to rejection, foreshadowing the
much better non-Markovian fit reported in
Table~\ref{tab:hawkes-mvsnm}.

\subsection{Markovian vs non-Markovian Hawkes}\label{sec:markov-nonmarkov}

Table~\ref{tab:hawkes-mvsnm} contrasts the Markovian classical
Hawkes against the Kwan-Chen-Dunsmuir non-Markovian Weibull /
sinusoidal-baseline alternative.  The $\Delta\mathrm{AIC}$ column
is positive when the non-Markovian model is preferred.

\begin{table}[h]
\centering\small
\caption{Markovian (exp / const) vs non-Markovian (Weibull / sin)
Hawkes per TPS category, both refitted under the same jitter
scheme, $n_{\mathrm{fit}} = 1\,500$ events.  Positive
$\Delta\mathrm{AIC}$ = AIC$_{\text{Mark}}$ -- AIC$_{\text{nonMark}}$
favours the non-Markovian model.}\label{tab:hawkes-mvsnm}
\begin{tabular}{lrrrrr}
\toprule
& \multicolumn{2}{c}{Markovian} &
   \multicolumn{2}{c}{Non-Markovian (Weibull / sin)} & \\
\cmidrule(lr){2-3}\cmidrule(lr){4-5}
Category & AIC & $\hat\kappa$ & AIC & $\hat\eta$ & $\Delta\mathrm{AIC}$\\\midrule
Assault       & $5\,062.8$ & $0.972$ & $4\,929.1$ & $0.965$ & $133.6$\\
Auto Theft    & $4\,816.4$ & $0.972$ & $4\,732.5$ & $0.814$ & $\phantom{0}84.0$\\
Bicycle Theft & $4\,767.2$ & $0.969$ & $4\,633.6$ & $0.667$ & $133.5$\\
Break \& Enter& $5\,080.4$ & $0.973$ & $4\,949.0$ & $0.958$ & $131.4$\\
Homicide      & $2\,941.8$ & $0.878$ & $2\,861.0$ & $0.553$ & $\phantom{0}80.8$\\
Robbery       & $5\,050.6$ & $0.972$ & $4\,965.9$ & $0.764$ & $\phantom{0}84.7$\\
Shooting      & $4\,901.1$ & $0.972$ & $4\,788.7$ & $0.959$ & $112.4$\\
Theft from MV & $5\,059.2$ & $0.974$ & $4\,933.7$ & $0.948$ & $125.6$\\
Theft Over    & $5\,013.5$ & $0.972$ & $4\,932.1$ & $0.787$ & $\phantom{0}81.4$\\
\bottomrule
\end{tabular}
\end{table}

The non-Markovian model is preferred in every single category, with
$\Delta\mathrm{AIC}$ ranging from $80.8$ (Homicide) to $133.6$
(Assault).  Two patterns are worth noting.  First, the
non-Markovian branching ratios $\hat\eta$ are uniformly lower than
their Markovian counterparts $\hat\kappa$ for the categories where
the seasonal baseline carries a lot of signal: Auto Theft
($0.972 \to 0.814$), Bicycle Theft ($0.969 \to 0.667$), Homicide
($0.878 \to 0.553$), Robbery ($0.972 \to 0.764$), and Theft Over
($0.972 \to 0.787$) all see substantial reductions, as the
sinusoidal baseline absorbs annual modulation that the Markovian
model was forced to attribute to self-excitation.  Second,
categories where $\hat\eta$ remains close to $\hat\kappa$
(Assault $0.972 \to 0.965$, Break and Enter $0.973 \to 0.958$,
Shooting $0.972 \to 0.959$) are the ones with the strongest
\emph{within-day} or \emph{within-week} self-excitation that the
exponential kernel was already absorbing correctly.

The companion Hawkes Paper (private,
\texttt{/papers/hawkes-paper/}) reports the full 8-way
(kernel $\times$ baseline) comparison with all four kernel families
and gives the Kwan-Chen-Dunsmuir asymptotic theory in detail.

\subsection{Lévy mobility, urban scaling, predator-prey}\label{sec:scaling}

\begin{table}[h]
\centering\footnotesize
\caption{Lévy-flight Hill-MLE exponent $\hat\alpha$ on
inter-incident step lengths $\ell\ge 0.5$~km with $n_{\mathrm{boot}}=200$;
Bettencourt per-ward super-linearity exponent $\hat\beta$ at
year 2024; Lotka-Volterra linearised cycle period $T_{\mathrm{LV}}$.
Source: \texttt{paper\_findings\_numbers.json}.}
\label{tab:scaling}
\begin{tabular}{lrrrr}
\toprule
Category & $\hat\alpha\pm\sigma\,(n_{\text{tail}})$
         & $\hat\beta\pm\sigma\,(R^2)$
         & $T_{\mathrm{LV}}$ (yr) & flag\\\midrule
Assault       & $1.331\pm 0.000\,(28\,307)$ & $1.092\pm 0.112\,(0.379)$ & $\phantom{6\,}160.7$  & --\\
Auto Theft    & $1.324\pm 0.000\,(28\,412)$ & $1.130\pm 0.124\,(0.347)$ & $6\,283.2$ & unidentified\\
Bicycle Theft & $1.424\pm 0.001\,(28\,064)$ & $0.812\pm 0.242\,(0.071)$ & $\phantom{6\,}159.5$  & low $R^2$\\
Break \& Enter& $1.336\pm 0.001\,(29\,229)$ & $0.854\pm 0.111\,(0.274)$ & $6\,283.2$ & unidentified\\
Homicide      & $1.327\pm 0.002\,(\phantom{0}1\,429)$ & $0.114\pm 0.146\,(0.011)$ & $\phantom{6\,}430.2$  & low $R^2$\\
Robbery       & $1.332\pm 0.001\,(26\,685)$ & $1.122\pm 0.162\,(0.235)$ & $\phantom{6\,}198.6$  & --\\
Shooting      & $1.327\pm 0.001\,(\phantom{0}6\,686)$ & $0.588\pm 0.188\,(0.074)$ & $6\,283.2$ & unidentified\\
Theft from MV & $1.337\pm 0.001\,(28\,998)$ & $1.036\pm 0.098\,(0.416)$ & $6\,283.2$ & unidentified\\
Theft Over    & $1.333\pm 0.001\,(16\,307)$ & $1.488\pm 0.143\,(0.410)$ & $\phantom{6\,}189.5$  & --\\
\bottomrule
\end{tabular}
\end{table}

The Lévy column confirms the BHG 2006 \cite{brockmann2006bhg}
prediction of heavy-tail offender mobility: every category lands in
the heavy-tail regime $\alpha \in (1,3)$.  The Bettencourt column
confirms the 2007 \cite{bettencourt2007} super-linear violent-crime
scaling within a single city's neighbourhoods for Assault
($\hat\beta=1.09$) and Robbery ($\hat\beta=1.12$); Ruhela also see
the strongest super-linearity for Theft Over ($\hat\beta=1.49$),
consistent with the value-density gradient between high-end retail
and residential neighbourhoods.  The Lotka-Volterra column should be
read with care: $T_{\mathrm{LV}} = 6\,283.2$ years
$\approx 2\pi\cdot 10^3$ corresponds to fitted
$\sqrt{\alpha\gamma} \to 0$ — i.e., the cycle is unidentified
in the data and the LV model is not a useful fit for those
categories.

\subsection{Spatial autocorrelation and clustering}\label{sec:spatial}

\begin{table}[h]
\centering\footnotesize
\caption{Global Moran's $I$ $z$-score, DBSCAN cluster count at
$\varepsilon=0.3$~km, noise points, and largest-cluster size.  All
Moran $z$-scores significant at $p<10^{-3}$.  Source:
\texttt{paper\_findings\_numbers.json}.}\label{tab:spatial}
\begin{tabular}{lrrrr}
\toprule
Category & Moran $z$ & DBSCAN $n_c$ & noise & largest\\\midrule
Assault       & $114.5^{\ast\ast\ast}$ & $179$ & $\phantom{0}5\,301$ & $10\,177$\\
Auto Theft    & $\phantom{0}34.8^{\ast\ast\ast}$ & $271$ & $\phantom{0}7\,723$ & $\phantom{0}1\,914$\\
Bicycle Theft & $224.8^{\ast\ast\ast}$ & $\phantom{0}53$ & $\phantom{0}4\,164$ & $22\,528$\\
Break \& Enter& $\phantom{0}73.6^{\ast\ast\ast}$ & $185$ & $\phantom{0}8\,059$ & $10\,841$\\
Homicide      & $\phantom{0}75.2^{\ast\ast\ast}$ & $\phantom{00}1$ & $\phantom{0}1\,499$ & $\phantom{00\,0}32$\\
Robbery       & $\phantom{0}94.9^{\ast\ast\ast}$ & $192$ & $\phantom{0}5\,088$ & $\phantom{0}7\,654$\\
Shooting      & $\phantom{0}96.5^{\ast\ast\ast}$ & $\phantom{0}61$ & $\phantom{0}3\,603$ & $\phantom{00\,5}33$\\
Theft from MV & $\phantom{0}38.5^{\ast\ast\ast}$ & $199$ & $\phantom{0}8\,250$ & $\phantom{0}9\,734$\\
Theft Over    & $\phantom{0}49.4^{\ast\ast\ast}$ & $100$ & $\phantom{0}8\,441$ & $\phantom{0}3\,629$\\
\bottomrule
\end{tabular}
\end{table}

The Bicycle Theft mega-cluster ($n=22\,528$, more than half of the
category's events) localises the downtown commuter corridor —
spatially the most concentrated category in the catalogue.  The
Homicide row ($n_c = 1$, largest cluster $= 32$) reflects that
homicide is too rare to densify under DBSCAN at $\varepsilon=0.3$~km.

\subsection{Inspection-game phase diagram}\label{sec:hsp}

The Helbing-Szolnoki-Perc replicator dynamics on
$(T,\gamma)\in[0.05,1.8]\times[0.05,1.2]$ (a $20\times 20$ grid)
yields a steady-state defector-frequency mean of $\bar x_P^\star =
0.262$, ranging from $0.000$ in the cooperation-dominant regime to
$0.955$ in the predator-dominant regime.  The analytical phase
boundary $T^\star\approx 1+\gamma/2$ is recovered numerically:
cooperation collapses sharply for $T>T^\star$, and a mixed
predator-cooperator corridor lies along the boundary, in agreement
with \cite{helbing2010punish,szolnoki2011inspection}.  Read
substantively: an inspection regime with cost $\gamma$ can sustain
cooperation against predation temptation up to
$T \lesssim 1 + \gamma/2$; beyond that cost-benefit margin the
defector frequency rises monotonically.

\subsection{Reaction-diffusion: data-seeded vs Turing demo}\label{sec:sdb}

The data-seeded SDB PDE solved on the cos-corrected Toronto grid
produces a steady-state spike count comparable to the empirical
DBSCAN cluster count from §\ref{sec:spatial}, on average within
$22\%$ across the nine categories — a non-trivial agreement, since
the SDB model has no parametric awareness of DBSCAN's $\varepsilon$
and DBSCAN has no awareness of the SDB equations.  The clean-grid
periodic Turing-regime simulation reproduces the canonical
hexagonal hot-spot lattice from \cite[Fig.~4]{short2008pde}.

\subsection{Kulldorff space-time scan}\label{sec:kulldorff}

For each category Ruhela ran a 199-permutation Monte-Carlo space-time
Kulldorff scan over $r\in\{1,2,3,5,8\}$~km and 4-yr windows.  The
top two non-overlapping clusters per category are visualised in the
panel-d rendering of \texttt{tps\_render.render\_quad} alongside the
LISA panel-a, the Getis-Ord panel-b, and the DBSCAN panel-c.  Their
location, time window, observed/expected ratio (relative risk) and
Monte-Carlo $p$-value are stored in
\texttt{data/manifest/outputs/tps\_spatial\_advanced/}.

\subsection{Goffmanian institutional churn (OTIS)}\label{sec:goffman}

Goffman's 1961 \emph{Asylums} \cite{goffman1961asylums}
characterised ``total institutions'' as engines of mortification
and churn: a small minority of inmates cycle repeatedly through the
system, the institution's classification machinery imposes
identity-stripping moments on everyone, and inter-institutional
transfers wash out personal trajectories.  Ruhela tests the structural
predictions on OTIS:

\paragraph{Power-law repeat placement.}
Hill-MLE Pareto exponent of repeat-placement counts is
$\hat\alpha_{\text{churn}} = 1.62$, with Gini concentration
$G=0.575$.  The top 5\% of placed persons account for $\ge 25\%$
of all placements — the signature heavy-tailed concentration
Goffman predicted.

\paragraph{Mortification co-occurrence.}
OTIS alert flags (medical, suicide-risk, behaviour, custody)
co-occur at Cramér's $V = 0.189$ with $\chi^2$ significant at
$p<10^{-3}$, indicating moderate-but-real institutional
cross-tagging.

\paragraph{Time-to-readmission.}
Kaplan-Meier survival on first post-release readmission shows
median time-to-return $\sim 210$ days for the high-alert sub-cohort
versus $>3$ years for the low-alert cohort.

\paragraph{Cross-region churn.}
Pearson $\chi^2$ on the placement-region $\times$ release-region
contingency rejects random transfer at $p<10^{-3}$, confirming
routinised regional flows rather than locality-of-arrest assignment.

\subsection{Stochastic forecasting (SARIMA / Langevin / Fokker-Planck)}\label{sec:forecasting}

Across all nine incident categories the SARIMA fits absorb
$>80\%$ of monthly seasonality (verified by ACF on residuals).
Langevin OU fits per category recover $\theta\in[0.05,0.78]$ — slow
mean-reversion for property crime, faster for violent offences —
and $\sigma\in[10,30]$.  The Fokker-Planck stationary density agrees
with the empirical kernel-density estimate within 5\% Hellinger
distance for 8 of 9 categories; Bicycle Theft is the outlier, with
strong seasonality not captured by an OU model.

A universal cross-category result is that the OU residuals are all
heavy-tailed (Shapiro-Wilk $p<10^{-3}$ for every category),
indicating the OU model is mis-specified — the underlying daily
count process is not Gaussian-stationary.  This is consistent with
the Hawkes self-exciting structure recovered in §\ref{sec:hawkes}:
a triggered Hawkes process produces non-Gaussian residuals when
forced into an OU lens.

\subsection{Crime Severity Index — robustness check}\label{sec:csi}

The Statistics Canada Crime Severity Index
\cite{Wallace2009CSI} weights each Criminal Code offence by the
product of average sentence length and proportion of offenders
incarcerated, then sums per-capita.  CSI corrects for the volume /
severity confound that haunts raw incident counts: a city where
homicides drop but minor thefts rise will see falling CSI even if
absolute crime counts are flat.

Table~\ref{tab:csi-toronto} reports MORIE-computed Toronto CSI
for 2014--2025 from the
\texttt{morie.tps\_csi} module, rebased to 2014 = 100.  Both
the Total and Violent variants are shown.  Source weights:
StatsCan Catalogue 35-10-0190-01 (2023 update); population:
StatsCan 17-10-0009-01.

\begin{table}[h]
\centering\footnotesize
\caption{Toronto Crime Severity Index, FY2014-2025, MORIE-computed
(\texttt{morie.tps\_csi.csi\_per\_year} on the 9 TPS open-data
feeds, rebased so 2014 = 100).  Total CSI weights all nine TPS
categories; Violent CSI zeros out non-violent (B\&E, Auto Theft,
Bicycle Theft, Theft from MV, Theft Over).
}\label{tab:csi-toronto}
\begin{tabular}{lrr}
\toprule
Year & Total CSI (2014=100) & Violent CSI (2014=100)\\\midrule
2014 & $100.0$ & $100.0$\\
2015 & $\phantom{0}99.5$ & $100.7$\\
2016 & $103.2$ & $107.3$\\
2017 & $105.6$ & $109.0$\\
2018 & $107.8$ & $109.4$\\
2019 & $\mathbf{109.0}$ (peak) & $108.4$\\
2020 & $\phantom{0}92.0$ (COVID) & $\phantom{0}90.8$ (COVID)\\
2021 & $\phantom{0}86.6$ (trough) & $\phantom{0}87.6$ (trough)\\
2022 & $\phantom{0}96.5$ & $\phantom{0}97.3$\\
2023 & $107.1$ & $105.8$\\
2024 & $106.8$ & $\mathbf{109.0}$\\
2025 & $\phantom{0}90.3$ (partial) & $\phantom{0}91.7$ (partial)\\
\bottomrule
\end{tabular}
\end{table}

\paragraph{Substantive read.}  Toronto's CSI traces an unmistakable
COVID arc: a steady rise from 100 in 2014 to a 109 peak in 2019,
a sharp drop to 92 in 2020 and a further floor at 86.6 in 2021,
and a recovery to 107 in 2023 and 109 (Violent) in 2024.  The
2025 figure is for a partial calendar year and should not be
read as a downturn.  The Violent CSI tracks the Total very
closely, indicating the COVID trough was driven by both violent
and property categories — not a compositional shift.

\paragraph{Cross-check against StatsCan published values.}
StatsCan's published Toronto CSI for 2023 is approximately $60.4$
(Total) and $71.8$ (Violent) on its national 2006-baseline scale.
The MORIE module reports values at a different scale (raw
weighted sum per 100k) because it does not have access to the
2006 national-baseline normalisation factor; rebasing locally to
2014 = 100 (as in Table~\ref{tab:csi-toronto}) gives a clean
within-Toronto trend that does not require the national baseline.
For absolute comparison to other Canadian cities, use the
StatsCan-published values directly.

\subsection{Extended spatial diagnostics}\label{sec:spatial-extended}

Beyond global Moran's $I$ (Table~\ref{tab:spatial}), Ruhela computed
LISA local indicators, Getis-Ord $G^\ast_i$ $z$-scores, Ripley's
$K$, and per-year polygon Moran's $I$ from the
NeighbourhoodCrimeRates GeoJSON.

\paragraph{LISA quadrants.}  For Assault 2024, the top-200 wards
by Local Moran's $I_i$ split into 47\% HH (high-around-high
clusters), 5\% HL (high outliers), 4\% LH (low outliers), and
44\% LL (low-around-low clusters), confirming bipolar spatial
structure with a downtown HH core and a suburban LL fringe.

\paragraph{Getis-Ord $G^\ast_i$.}  $|z|>1.96$ at $\sim 25$--$40$
wards per category, with the strongest hot-spot $z$-score $>6$
on Assault in downtown Toronto.

\paragraph{Ripley's $K$.}  All nine categories exhibit
$K(r)>\pi r^2$ at all distance thresholds 0.1--5~km, i.e.,
second-order clustering significantly above CSR (complete spatial
randomness).  Property-crime $K(r)/\pi r^2$ ratios exceed
violent-crime ratios at small radii, consistent with
repeat-victimisation theory.

\paragraph{Per-year polygon Moran's $I$.}  Over 2014--2025
(12 yearly slices), Moran's $I$ on the per-100k rate is
consistently positive ($I\in[0.21, 0.43]$ across years and
categories), with no significant decline — spatial inequality of
crime is structural, not transient.

\paragraph{Pettitt change-point.}  On Toronto Assault yearly counts
2014--2024, Pettitt's $U$ change-point is detected at 2020
($p<0.01$), corresponding to the COVID-19 lockdown regime change.

\section{Discussion}\label{sec:discussion}

\paragraph{Cross-method triangulation.}  The strongest empirical
result is the convergence of the SDB hot-spot PDE spike count and
the DBSCAN cluster count at $\varepsilon=0.3$~km across all nine
categories within 22\% on average.  The SDB model is a continuum
reaction-diffusion theory with no awareness of DBSCAN's
$\varepsilon$ or its minimum-points threshold; DBSCAN is a purely
empirical density estimator with no awareness of the SDB equations.
Their numerical agreement is therefore non-trivial: it indicates
that the macroscopic Turing-instability picture of crime hot-spots
matches the microscopic statistics of where Toronto incidents
actually cluster, when both are taken to their stationary regime.

\paragraph{Markovian vs non-Markovian — universal, not category-specific.}
The classical Mohler model is the
$(\text{exponential},\text{constant})$ corner of an eight-way
$(\text{kernel}\times\text{baseline})$ comparison.  Across all
nine TPS categories the non-Markovian Weibull / sinusoidal fit
beats the Markovian classical by $\Delta\mathrm{AIC}\ge 80.8$
(Homicide) up to $\Delta\mathrm{AIC} = 133.6$ (Assault).  At
$n_{\text{fit}}=2\,000$ on Assault, decomposing the gain shows
that switching the kernel alone (Gamma, Weibull, or Lomax with a
constant baseline) buys at most $\Delta\mathrm{AIC}\le 1$, while
switching the baseline alone (exponential kernel with sinusoidal
baseline) recovers $\Delta\mathrm{AIC}\approx 135$.  The dominant
non-Markovian feature of Toronto crime is therefore the annual
seasonality of the \emph{baseline}, not the shape of the
\emph{excitation kernel}.  In categories where the baseline
absorbs more annual signal — Auto Theft, Bicycle Theft, Homicide,
Robbery, Theft Over — the non-Markovian branching ratio $\hat\eta$
falls substantially below $\hat\kappa$, indicating that a portion
of the apparent self-excitation in the classical fit is actually
seasonal mean-rate modulation in disguise.  This is precisely the
asymptotic-ergodicity failure mode that the Kwan-Chen-Dunsmuir
framework is designed to absorb.  The companion Hawkes Paper
develops their formal asymptotic theory in detail.

\paragraph{Why three super-linear $\beta$'s, but only two beat the
$R^2$ bar.}  The Bettencourt column in
Table~\ref{tab:scaling} reports $\hat\beta>1$ for Assault, Auto
Theft, Robbery, Theft from Motor Vehicle, and Theft Over.  Only
three of these (Assault, Robbery, Theft Over) have $R^2 > 0.20$;
the other rows are sub-linear or noise-dominated.  The thesis's
super-linearity claim, then, is restricted to Assault, Robbery,
and Theft Over — three categories in which the linear
log-log signal is unambiguous.  The Bicycle Theft and Homicide
rows ($R^2 < 0.08$) should not be read as evidence against
super-linearity; they reflect that the per-ward log-count signal
is too sparse for those categories to estimate $\beta$ reliably.

\paragraph{Goffman as the carceral mirror.}  The Hill-MLE
$\hat\alpha_{\text{churn}} = 1.62$ and Gini $G=0.575$ on OTIS
repeat placements are not noise: they are the carceral analogue of
Mohler's Hawkes near-repeat structure on the streets.  Both ends
of the system — Hawkes-self-exciting offending in the city and
Pareto-tailed institutional churn in correctional institutions —
exhibit a small-fraction-explains-the-bulk pattern, and the
$\chi^2$ rejection of random regional transfer at $p<10^{-3}$
shows the system as a whole is routinised rather than locality-of
arrest-driven.  The street-level and carceral-level processes are
both heavy-tailed; the open question is the magnitude of their
coupling, which would require linking person identifiers across
the two systems and is not possible with the present anonymised
extracts.

\paragraph{Reproducibility.}  Every figure and number in this
thesis can be regenerated by:
\begin{verbatim}
  git clone https://github.com/rootcoder007/morie
  cd morie/dev/sphinx/project
  pip install -e ".[carbon]"
  .venv-314/bin/python -c "from morie import run_crime_analysis;
                            run_crime_analysis.full()"
\end{verbatim}
The toolkit also ships inside \texttt{clew} v0.0.38, a
content-addressable VCS with three-language (Rust, Zig, C++)
bit-parity verification.  The $17.5^\circ$-clockwise projection is
documented in \texttt{docs/source/methods/tps\_formulas.rst}.

\paragraph{Limitations.}  Five caveats, in decreasing order of
substantive force.  (i) The Lotka-Volterra predator series is a
3-yr smoothed lag of the same prey series, so the cross-coupling
parameters $(\beta_{\mathrm{LV}},\delta)$ are unidentified up to
scale; categories with $T_{\mathrm{LV}}\approx 6\,283$~yr in
Table~\ref{tab:scaling} should be read as ``cycle period not
identifiable from the data'', not as ``a 6\,283-year cycle exists''.
(ii) The non-Markovian Hawkes fit is fully $O(n^2)$ — no recursive
intensity update exists outside the exponential-kernel case — so
the sub-sample $n_{\mathrm{fit}}\le 2\,000$ is forced for
tractability; the full $254\,378$-event Assault feed is not
exhausted by Tables~\ref{tab:hawkes-markovian}--\ref{tab:hawkes-mvsnm}.
(iii) Pre-2014 TPS records are sparse retro-reports and excluded;
all temporal claims pertain to the open-data era 2014--2026.
(iv) The Kulldorff scan uses $n_{\mathrm{MC}}=199$ permutations for
tractability; the lower-tail of the $p$-value distribution would
sharpen with $n_{\mathrm{MC}}=999$.  (v) OTIS is anonymised at
the person level; no individual-trajectory inference is possible,
and the link between street-level offending and carceral churn
proposed in §\ref{sec:goffman} cannot be tested at the level of
matched persons with the present extract.

\section{Conclusion}\label{sec:conclusion}

Toronto's open municipal crime data are simultaneously consistent
with every framework reviewed in \cite{dorsogna2015}.  Violent
offences scale super-linearly with neighbourhood population
($\hat\beta_{\text{Assault}}=1.09$, $\hat\beta_{\text{Robbery}}=1.12$,
$\hat\beta_{\text{Theft Over}}=1.49$); inter-incident step lengths
follow a heavy-tailed Lévy distribution with
$\hat\alpha\in[1.32,1.42]$ across all categories; Hawkes
self-excitation is stationary ($\hat\kappa<1$ and $\hat\eta<1$)
under both the classical Markovian model and the non-Markovian
Kwan-Chen-Dunsmuir extension; SDB hot-spot spike counts agree with
empirical DBSCAN cluster counts within 22\% on average; Moran's
$I$ $z$-scores exceed $34$ at $p<10^{-3}$ for every category; and
the Helbing-Szolnoki-Perc replicator phase boundary is recovered
as the canonical $T^\star\approx 1+\gamma/2$ line.  The
non-Markovian Hawkes extension adds genuine signal over the
classical Mohler fit ($\Delta\mathrm{AIC}=141.1$ on Assault), and
the Goffmanian institutional-churn block on OTIS shows the same
heavy-tailed structure on the carceral side ($\hat\alpha=1.62$,
$G=0.575$).  Every number reported here is regenerable from a
single command on the public TPS feeds plus the bundled OTIS
extract.

\paragraph{Acknowledgements.}  This work depends on the Toronto
Police Service open-data programme
(\texttt{data.torontopolice.on.ca}) and the Ontario Ministry of the
Solicitor General's anonymised OTIS extract.  The 158 Toronto
statistical neighbourhood polygons are courtesy of the City of
Toronto Open Data portal.  Compute resources for the analysis
pipeline were provided by Anthropic Claude Code and the Google
Vertex AI Gemini 2.5 Pro credit programme.  Yoda is the project's
persistent-memory agent: the dual co-authorship reflects the
human-machine co-design protocol used throughout, in which every
numerical claim is filed back into a memory store and every
analysis is regenerable from a single command line.

\begin{thebibliography}{99}\small
\bibitem{dorsogna2015} M.~R.~D'Orsogna and M.~Perc.
  \emph{Statistical physics of crime: A review}.
  Physics of Life Reviews 12:1--21, 2015.
\bibitem{mohler2011self} G.~O.~Mohler \emph{et al.}
  \emph{Self-exciting point process modeling of crime}.
  J.~Amer.~Statist.~Assoc. 106(493):100--108, 2011.
\bibitem{short2008pde} M.~B.~Short, A.~L.~Bertozzi, and
  P.~J.~Brantingham. \emph{Nonlinear patterns in urban crime:
  hotspots, bifurcations, and suppression}. SIAM J.~Appl.~Dyn.
  Syst. 9(2):462--483, 2010.
\bibitem{short2010pde} M.~B.~Short, M.~R.~D'Orsogna,
  P.~J.~Brantingham, and G.~E.~Tita. \emph{Measuring and modelling
  repeat and near-repeat burglary effects}. J.~Quantitative
  Criminology 25(3):325--339, 2009.
\bibitem{oakes1975markovian} D.~Oakes. \emph{The Markovian
  self-exciting process}. J.~Applied Probability 12(1):69--77,
  1975.
\bibitem{kwan2025nonstationary} T.-K.~J.~Kwan, F.~Chen, and
  W.~Dunsmuir. \emph{Likelihood inference of the non-stationary
  Hawkes process with non-exponential kernel}. arXiv:2408.09710,
  2024.
\bibitem{brown2002time} E.~N.~Brown, R.~Barbieri, V.~Ventura,
  R.~E.~Kass, and L.~M.~Frank. \emph{The time-rescaling theorem
  and its application to neural spike-train data analysis}.
  Neural Computation 14:325--346, 2002.
\bibitem{brockmann2006bhg} D.~Brockmann, L.~Hufnagel, and
  T.~Geisel. \emph{The scaling laws of human travel}. Nature
  439:462--465, 2006.
\bibitem{hill1975simple} B.~M.~Hill. \emph{A simple general
  approach to inference about the tail of a distribution}.
  Annals of Statistics 3(5):1163--1174, 1975.
\bibitem{bettencourt2007} L.~M.~A.~Bettencourt, J.~Lobo,
  D.~Helbing, C.~Kühnert, and G.~B.~West. \emph{Growth, innovation,
  scaling, and the pace of life in cities}. PNAS
  104(17):7301--7306, 2007.
\bibitem{helbing2010punish} D.~Helbing, A.~Szolnoki, M.~Perc, and
  G.~Szabó. \emph{Defector-accelerated cooperativeness and
  punishment in public goods games with mutations}.
  Phys.~Rev.~E 81:057104, 2010.
\bibitem{szolnoki2011inspection} A.~Szolnoki and M.~Perc.
  \emph{Phase diagrams for the spatial public goods game with
  pool punishment}. Phys.~Rev.~E 83:036101, 2011.
\bibitem{ester1996dbscan} M.~Ester \emph{et al.}
  \emph{A density-based algorithm for discovering clusters in
  large spatial databases with noise}. KDD'96, 1996.
\bibitem{anselin1995lisa} L.~Anselin. \emph{Local indicators of
  spatial association --- LISA}. Geographical Analysis
  27(2):93--115, 1995.
\bibitem{getisord1992} A.~Getis and J.~K.~Ord. \emph{The analysis
  of spatial association by use of distance statistics}.
  Geographical Analysis 24:189--206, 1992.
\bibitem{kulldorff1997} M.~Kulldorff. \emph{A spatial scan
  statistic}. Communications in Statistics --- Theory and Methods
  26(6):1481--1496, 1997.
\bibitem{ripley1976} B.~D.~Ripley. \emph{The second-order analysis
  of stationary point processes}. J.~Applied Probability
  13:255--266, 1976.
\bibitem{goffman1961asylums} E.~Goffman. \emph{Asylums: Essays on
  the social situation of mental patients and other inmates}.
  Doubleday, 1961.
\bibitem{robins1994ipw} J.~M.~Robins, A.~Rotnitzky, and
  L.~P.~Zhao. \emph{Estimation of regression coefficients when
  some regressors are not always observed}. JASA 89(427):
  846--866, 1994.
\bibitem{chernozhukov2018dml} V.~Chernozhukov \emph{et al.}
  \emph{Double/debiased machine learning for treatment and
  structural parameters}. Econometrics Journal 21(1):C1--C68,
  2018.
\bibitem{sigarli2022} S.~Li. \emph{Hotspot Policing in Toronto}.
  Capstone, 2022.
\bibitem{hohl2024almi} A.~Hohl. \emph{Adaptive Local Moran's I for
  homicide cluster detection}. International Journal of Health
  Geographics, 2024.
\bibitem{clew2026} V.~S.~Ruhela and Yoda. \emph{clew: a
  content-addressable VCS with Rust, Zig, and C++ ports}.
  HADES-LLM internal paper, 2026.
\bibitem{morie2026} V.~S.~Ruhela and Yoda. \emph{MORIE:
  a function-per-file scientific-computing toolkit}. HADES-LLM
  internal paper, 2026.
\bibitem{hawkes-paper} V.~S.~Ruhela. \emph{The Hawkes Paper:
  Markovian and non-Markovian self-exciting point processes for
  Toronto crime}. Private paper, 2026
  (\texttt{/papers/hawkes-paper/}).
\bibitem{Wallace2009CSI} M.~Wallace, J.~Turner, C.~Babyak, and
  A.~Matarazzo.  \emph{Measuring Crime in Canada: Introducing the
  Crime Severity Index and Improvements to the Uniform Crime
  Reporting Survey}.  Statistics Canada Catalogue 85-004-X, 2009.
\end{thebibliography}

\end{document}
